\section{Sibling EICrecon Integration and Modes}
The public integration is designed for a compatible sibling EICrecon checkout. The integration guide specifies the additions to copy, the two build-registration changes to apply, and the runtime variables required by the local worker and reconstruction wrappers. No modification to a sibling checkout is performed by this report build.

\begin{table}[htbp]\centering\small
\caption{Selection semantics.}\begin{tabular}{L{0.24\linewidth}L{0.60\linewidth}}\toprule
Mode & Behavior\\\midrule
Classical & C++ applies the deterministic first finite minimum.\\
\texttt{qiskit\_shadow} & A local selection is requested and recorded, while C++ applies the classical action.\\
\texttt{qiskit\_active} & C++ applies a local index only after strict validation and exact-wire digest agreement.\\\bottomrule
\end{tabular}\end{table}

Shadow mode separates observation from action: it records what the local selector proposed while preserving the classical clustering transition. Active mode changes that final selection step only when the reply is valid. The bare EICrecon configuration defaults \texttt{quantumFailClosed=false} and may classically fall back. The public active wrapper forces \texttt{true}; the public shadow wrapper defaults to \texttt{true} and offers only an exploratory \texttt{false} override. Oversize, timeout, and invalid-reply behavior depends on that configuration. An exact-zero pair distance is applied locally by C++ without a worker request or digest.

\measured The integration retains C++ ownership throughout both modes. The worker does not construct event data, mutate jets, or write EDM output. This division is essential for interpreting a selector comparison: it is a comparison of candidate-choice behavior within a fixed reconstruction loop, not a replacement of the reconstruction framework.
\clearpage
